\documentclass[11pt]{article}
\usepackage[margin=1in]{geometry}
\usepackage{amsmath,amssymb}
\usepackage{booktabs}
\usepackage{hyperref}
\usepackage{url}
\usepackage{listings}
\lstset{basicstyle=\ttfamily\small,breaklines=true}

\title{Epistemic Shadow Lanes: Compiler-Generated Uncertainty Propagation for GPU Kernels}
\author{Demetrios C. Agourakis\\
Biomaterials and Regenerative Medicine Post-Graduate Program,\\
Pontif\'{i}cia Universidade Cat\'{o}lica de S\~{a}o Paulo (PUC-SP), Brazil\\
Faculdade S\~{a}o Leopoldo Mandic, Campinas, Brazil\\
ORCID: 0009-0001-8671-8878}
\date{July 2026}

\begin{document}
\maketitle
\begin{abstract}
We present a compiler synthesis that couples first-order GUM-style (JCGM 100:2008) uncertainty propagation with GPU kernel emission, including WMMA tensor-core paths. The contribution is architectural integration --- not new uncertainty formulas --- but automatic shadow-register emission at the PTX level, typed via \texttt{Knowledge}$\langle T \rangle$, with optional correlated-input and second-order extensions. For each value register, the compiler maintains variance, validity, and provenance shadows that implement the law of propagation of uncertainty without programmer intervention. We study whether epistemic metadata can guide five GPU-specific execution-time optimizations under a formal uncertainty budget: (1) warp-level ballot voting on validity predicates for budget-bounded dual-path execution, (2) Shannon entropy of the relative-uncertainty distribution for kernel variant selection, (3) provenance-aware DAG scheduling with automatic stream parallelism, (4) speculative guards for refinement kernels below an uncertainty threshold, and (5) provenance-weighted kernel fusion scoring. A pharmacokinetic case study with four noisy sensors eliminates approximately 190 lines of manual variance mathematics per CUDA kernel. Instruction-count overhead is modelled exactly in Section 3.3; preliminary NVIDIA L4 GEMM measurements show a 2.19$\times$ geometric-mean latency ratio versus cuBLAS SGEMM before isolating shadow-register cost. We differentiate from Enzyme (Jacobian AD shadows), Puffin \cite{puffin2021} (CPU uncertainty compiler), and GBEES-GPU \cite{gbeesgpu2025} (grid-based Bayesian propagation), arguing that analytical GUM shadows through WMMA occupy an unaddressed design point.
\end{abstract}

\section{Introduction}

GPU-accelerated scientific computing routinely discards a critical property of its inputs: measurement uncertainty. A sensor reading of 50.0 ng/mL with standard uncertainty 2.0 ng/mL enters a GPU kernel as a bare \texttt{float}, and the output --- however many multiply-accumulate operations later --- is equally bare. The uncertainty is gone. The provenance is gone. The scientist must reconstruct both manually, if at all.

The Guide to the Expression of Uncertainty in Measurement (GUM, JCGM 100:2008) \cite{jcgm2008gum} standardizes how to propagate measurement uncertainty through computation. For a measurand $y = f(x_1, \ldots, x_N)$ with uncorrelated inputs, the combined standard uncertainty is:

\[u_c(y) = \sqrt{\sum_{i=1}^{N} \left(\frac{\partial f}{\partial x_i}\right)^2 u^2(x_i)}\]

This is a \emph{mechanical transformation} of the same arithmetic the programmer already writes. The partial derivatives (sensitivity coefficients) for elementary operations are known at compile time: for addition, $\partial(a+b)/\partial a = 1$; for multiplication, $\partial(ab)/\partial a = b$. This observation is the basis of our technique.

\subsection{The gap}

Five categories of prior work surround this space without occupying the same design point:

1. \textbf{Automatic differentiation on GPU.} Enzyme \cite{enzyme2021} inserts shadow registers alongside value registers in LLVM IR to propagate partial derivatives through GPU kernels. The structural mechanism --- compiler-generated shadow computation in GPU IR --- is analogous to our approach. However, Enzyme propagates \emph{Jacobian entries} ($\partial f/\partial x_i$), not \emph{combined variance} ($u_c^2(y)$). These are different mathematical objects with different composition rules (derivatives chain-multiply; variances quadrature-sum).

2. \textbf{Uncertain types.} Bornholt et al. \cite{uncertaint2014} introduced Uncertain$\langle T \rangle$ as a first-order type for uncertain data in managed languages. The type-level concept is valuable, but Uncertain$\langle T \rangle$ uses Monte Carlo sampling at runtime (not analytical propagation) and has no GPU backend.

3. \textbf{CPU uncertainty libraries and compilers.} Python's \texttt{uncertainties} \cite{uncertainties}, Julia's \texttt{Measurements.jl} \cite{measurementsjl}, and Puffin \cite{puffin2021} propagate uncertainty on CPU --- the last as an automatic uncertainty compiler that rewrites source to call intrusive UQ routines. All are CPU-only, with reported overheads of 50–1,500$\times$ for library approaches [6,7], and none emit GPU PTX or WMMA shadow lanes.

4. \textbf{GPU distribution propagation.} GBEES-GPU \cite{gbeesgpu2025} propagates high-dimensional probability distributions on GPU via grid-based Bayesian estimation. It targets full distribution evolution (Eulerian phase-space grids), not per-instruction analytical GUM variance shadows fused with general kernel arithmetic.

5. \textbf{Tensor-core numerics.} Blanchard et al. \cite{blanchard2020}, Fasi et al. \cite{fasi2021}, and Khattak and Mikaitis \cite{khattak2025} characterize \emph{rounding error} and hardware non-IEEE behaviour of NVIDIA tensor cores. Our work propagates \emph{measurement uncertainty} (a runtime quantity attached to input data) through WMMA instructions --- a different quantity with different composition rules.

No prior system combines: (i) compiler-emitted first-order GUM variance shadows, (ii) PTX/WMMA instruction-level integration, and (iii) language-level \texttt{Knowledge}$\langle T \rangle$ typing with provenance gates. This paper studies that synthesis.

\subsection{Contributions}

\textbf{Primary (design synthesis).}

1. \textbf{Epistemic shadow lanes} (Section 3): a compiler pass that, for each arithmetic instruction on a value register \texttt{%r_val}, emits GUM variance propagation on \texttt{%r_eps}, validity conjunction on \texttt{%p_valid}, and provenance OR-merge on \texttt{%r_prov}. Provenance uses bitwise OR (monotonic union) --- not XOR --- ensuring idempotent accumulation ($\text{prov}(a) \mathbin{|} \text{prov}(a) = \text{prov}(a)$).

2. \textbf{WMMA shadow tiles} (Section 2): dual \texttt{mma.sync.aligned} paths --- one for values, one for uncertainty fragments --- with element-wise quadrature across the inner dimension $k$, extending GUM product rules to tensor-core GEMM tiles.

3. \textbf{Five GPU-specific epistemic optimizations} (Section 4): warp-vote budget-bounded fast paths, entropy-gated dispatch, provenance-aware multi-stream scheduling, speculative refinement guards, and provenance-weighted fusion scoring --- each gated by formal uncertainty budgets rather than ad hoc thresholds.

\textbf{Secondary (engineering).}

4. \textbf{Multi-backend emission} (Section 3.4): the same shadow-lane logic targets PTX, Metal, and SPIR-V backends from a single compiler IR.

5. \textbf{Programmer effort elimination} (Section 5.1): a pharmacokinetic four-sensor fusion case study requires zero additional lines versus ~190 lines of manual CUDA variance plumbing.

\textbf{Honest scope (what we do not claim).}

\begin{itemize}
\item We do not claim the first uncertainty compiler; Puffin \cite{puffin2021} already rewrites source for intrusive UQ on CPU.
\item We do not claim novel GUM formulas; the propagation rules follow JCGM 100:2008 \cite{jcgm2008gum}.
\item We do not claim state-of-the-art GPU throughput; L4 GEMM measurements (Section 5.2) are preliminary and do not yet isolate shadow-register overhead from compiler maturity effects.
\end{itemize}


\section{Background: GUM Propagation Rules}

The GUM \cite{jcgm2008gum} specifies variance propagation for elementary operations assuming uncorrelated inputs. We implement the following rules, applied by the compiler at each arithmetic instruction:

\begin{table}[t]\centering
\begin{tabular}{lll}
\toprule
Source operation & Shadow emission: $u^2(\text{result})$ & GUM reference \\
\midrule
\texttt{add.f64 %c, %a, %b} & $u^2(a) + u^2(b)$ & Eq. (10), uncorrelated \\
\texttt{sub.f64 %c, %a, %b} & $u^2(a) + u^2(b)$ & Eq. (10) \\
\texttt{mul.f64 %c, %a, %b} & $b^2 u^2(a) + a^2 u^2(b)$ & Eq. (10) \\
\texttt{div.f64 %c, %a, %b} & $u^2(a)/b^2 + a^2 u^2(b)/b^4$ & Eq. (10) \\
\texttt{sqrt.f64 %c, %a} & $u^2(a) / (4a)$ & Eq. (13) \\
\bottomrule
\end{tabular}
\end{table}

For tensor core operations (\texttt{mma.sync.aligned}), the compiler emits \textbf{two} WMMA instructions per accumulation step: one on value fragments $(A, B)$ producing $C$, and one on uncertainty fragments $(u^2(A), u^2(B))$ with the same tile geometry. The output-tile variance is computed element-wise by quadrature (root-sum-square) across the inner dimension $k$:

\[u^2(C_{ij}) = \sum_k \left[ B_{kj}^2 \cdot u^2(A_{ik}) + A_{ik}^2 \cdot u^2(B_{kj}) \right]\]

This is the GUM product rule applied per $(i,j,k)$ triple, then summed --- not a single WMMA on Jacobians. The shadow WMMA reuses the value tile's $k$-loop structure but operates on variance fragment operands in f64 (or f32 on backends without native f64).


\section{Epistemic Shadow Lanes}

\subsection{1 Register Layout}

For each value register \texttt{%r_val} in the GpuKernelIr, the compiler allocates three shadow registers:

\begin{lstlisting}[basicstyle=\ttfamily\small]
%r_val   — point estimate (the measurand)
%r_eps   — variance u²(val)
%p_valid — validity predicate (true if uncertainty is well-defined)
%r_prov  — 64-bit OR provenance bitset (monotonic union)
\end{lstlisting}

The shadow registers are allocated during the HLIR-to-GpuKernelIr lowering pass and propagated through all subsequent optimization passes (constant folding, dead code elimination, register allocation).

\subsection{2 PTX Emission Example}

For a source-level multiplication \texttt{let c = a * b}, the compiler emits:

\begin{lstlisting}[basicstyle=\ttfamily\small]
// Value lane
mul.f64  %r_c_val, %r_a_val, %r_b_val;

// Variance lane: u²(c) = b²·u²(a) + a²·u²(b)
mul.f64  %r_t1, %r_b_val, %r_b_val;
mul.f64  %r_t1, %r_t1, %r_a_eps;
mul.f64  %r_t2, %r_a_val, %r_a_val;
mul.f64  %r_t2, %r_t2, %r_b_eps;
add.f64  %r_c_eps, %r_t1, %r_t2;

// Validity conjunction
and.pred  %p_c_valid, %p_a_valid, %p_b_valid;

// Provenance merge (OR is associative, commutative, idempotent)
or.b64   %r_c_prov, %r_a_prov, %r_b_prov;
\end{lstlisting}

A single source multiplication becomes 8 PTX instructions (1 value + 5 variance + 1 validity + 1 provenance). The shadow computation adds 7 instructions per multiply, 3 per add, and 9 per divide --- yielding an exact instruction-count ratio of 4–10$\times$ depending on operation mix (Section 3.3). Any wall-clock reduction below this ratio remains unverified and architecture-dependent (Section 5.2).

\subsection{3 Overhead Model}

For a kernel with $A$ additions, $M$ multiplications, and $D$ divisions, the instruction count overhead is:

\[\text{overhead} = \frac{A \cdot 4 + M \cdot 8 + D \cdot 10}{A + M + D}\]

The coefficients 4, 8, and 10 are exact instruction counts per operation class (Section 3.2). For pure addition kernels the ratio is exactly 4; for pure division kernels exactly 10. Mapping this ratio directly to a 4–10$\times$ \emph{wall-clock} range assumes uniform per-instruction latency --- an assumption we have not verified. Shadow operations may reuse operands already in registers and benefit from warp-level parallelism, but the magnitude of any wall-clock reduction is architecture-, occupancy-, and register-pressure-dependent (Section 5.2, Section 7.3).

\subsection{4 Multi-Backend Emission}

The same shadow lane logic targets three backends:

\begin{table}[t]\centering
\begin{tabular}{llll}
\toprule
Backend & Shadow precision & Tensor core support & Provenance \\
\midrule
PTX (NVIDIA CUDA) & f64 native & WMMA m16n8k16 & 64-bit OR \\
Metal (Apple MSL) & f32 (no native f64) & N/A & 32-bit OR \\
SPIR-V (Vulkan) & f64 via extension & N/A & 64-bit OR \\
\bottomrule
\end{tabular}
\end{table}


\section{GPU-Specific Optimizations}

\subsection{1 Warp-Vote Epistemic Fast-Path}

\textbf{Observation.} In many scientific workloads, large regions of data are well-characterized (high confidence, low uncertainty). Only boundary or anomalous regions require full uncertainty tracking. If all 32 lanes in a warp have valid data with uncertainty below a threshold, the shadow computation can be performed with reduced precision within a formal uncertainty budget.

\textbf{Mechanism.} The compiler generates dual-path kernels:

\begin{lstlisting}[basicstyle=\ttfamily\small]
// Phase 1: ballot — are ALL lanes valid?
vote.sync.ballot.b32  %r_ballot, %p_valid, 0xFFFFFFFF;
setp.eq.u32           %p_all_valid, %r_ballot, 0xFFFFFFFF;

// Phase 2: ballot — are ALL uncertainties below threshold?
setp.lt.f64           %p_eps_ok, %r_eps, THRESHOLD;
vote.sync.ballot.b32  %r_eps_ballot, %p_eps_ok, 0xFFFFFFFF;
setp.eq.u32           %p_all_ok, %r_eps_ballot, 0xFFFFFFFF;

@%p_all_ok bra FAST_PATH;

FULL_PATH:
    // 8 instructions per multiply (value + 5 variance + validity + provenance)
    ...
    bra MERGE;

FAST_PATH:
    // Budget-bounded uncertainty propagation (operator-specific bounds)
    // add/sub: u²_out ≤ 2·max(u²_a, u²_b)
    // mul:     u²_out ≤ (M_a² + M_b²)·max(u²_a, u²_b)  [requires operand envelopes]
    // div:     u²_out ≤ u²_a/m_b² + M_a²·u²_b/m_b⁴      [requires denom lower bound]
    // Validity: AND of inputs (preserved)
    // Provenance: OR of inputs (preserved)
    ...

MERGE:
    // Reconvergence
\end{lstlisting}

\textbf{Semantic invariant.} The fast path does \emph{not} skip shadow computation --- small uncertainty is not zero uncertainty. Instead, it uses operator-specific conservative bounds derived from operand envelopes. For addition/subtraction, $u^2_{\text{out}} \leq 2 \max(u^2_a, u^2_b)$. For multiplication, conservative bounds additionally require runtime or compile-time envelopes on operand magnitudes ($|a| \leq M_a$, $|b| \leq M_b$), giving $u^2_{\text{out}} \leq (M_a^2 + M_b^2) \max(u^2_a, u^2_b)$. For division, the bound requires a lower bound on the denominator ($|b| \geq m_b > 0$). These bounds are \emph{sound with respect to a user-specified uncertainty budget} (Definition 1 below), not "GUM-compliant" in the metrological sense --- the GUM provides the propagation law; the fast path is an optimization atop it.

This costs approximately 30% of full shadow overhead (the \texttt{budget_overhead} parameter in the divergence cost model), yielding a predicted speedup range of 1.3–1.7x depending on divergence percentage. The ballot overhead (2 instructions per check point) is amortized over the loop body.

\textbf{Definition 1 (Uncertainty Budget).} For a kernel output $y$, define the budget gap as:

\[\Delta(y) = u^2_{\text{fast}}(y) - u^2_{\text{full}}(y)\]

The fast path is \emph{budget-sound} if $\Delta(y) \leq B$ for a user-specified absolute bound $B$, or equivalently $\Delta(y) / u^2_{\text{full}}(y) \leq \rho$ for a relative bound $\rho$. The budget is \emph{local} (per-instruction): each operator's conservative bound overapproximates the full GUM result by at most a bounded factor. Composition across a kernel of $L$ instructions accumulates at most $L \cdot B_{\text{local}}$ total budget gap; tighter analysis via interval arithmetic is future work.

\textbf{Novelty claim.} \texttt{vote.sync.ballot} is a standard CUDA primitive used for reductions, stream compaction, and divergence management. Using it to evaluate \emph{epistemic validity predicates} --- "are all lanes' uncertainty values below a threshold?" --- to gate \emph{budget-bounded dual-path kernel execution} is, to our knowledge, a novel application.

\subsection{2 Entropy-Gated Kernel Dispatch}

\textbf{Observation.} The distribution of uncertainty values in the input buffer carries information about which kernel variant to dispatch. Concentrated uncertainty (all values similar) suggests the fast path will dominate; spread uncertainty suggests full GUM propagation is needed.

\textbf{Mechanism.} Before kernel launch, the host:

1. Samples $k$ uncertainty values from the input buffer (default $k = 256$)
2. Histograms them into $B$ bins (default $B = 16$)
3. Computes Shannon entropy: $H(\epsilon) = -\sum_{i=1}^{B} p_i \log_2 p_i$

\begin{table}[t]\centering
\begin{tabular}{lll}
\toprule
$H(\epsilon)$ & Dispatch & Rationale \\
\midrule
$< 1.0$ bit & Fast kernel & Concentrated; warp-vote fast path will dominate \\
$1.0 - 3.0$ bits & Adaptive & Mixed; both paths will activate \\
$\geq 3.0$ bits & Full GUM & Spread; full propagation needed everywhere \\
\bottomrule
\end{tabular}
\end{table}

\textbf{Novelty claim.} Existing adaptive kernel dispatch systems (Stream-K++ \cite{streamkpp2024}, KernelFoundry \cite{kernelfoundry2026}, Triton autotuning \cite{triton2019}) select variants based on \emph{performance metrics} (throughput, latency) or \emph{problem shape} (matrix dimensions). Using Shannon entropy of the \emph{data uncertainty itself} for dispatch is, to our knowledge, novel.

\subsection{3 Provenance-Aware DAG Scheduling}

\textbf{Observation.} After conventional dependence analysis has ruled out read/write and explicit dataflow dependencies, kernels operating on data from \emph{disjoint provenance sources} are treated as candidates for concurrent scheduling on separate CUDA streams. Provenance disjointness is a \emph{secondary scheduling heuristic} for identifying concurrency opportunities among otherwise independent kernels --- it does not substitute for alias analysis or read/write-set analysis.

\textbf{Mechanism.} The compiler:

1. Computes a provenance tag per kernel from its workload classification (bit-packed)
2. Builds a dependency graph: edge from $K_i$ to $K_j$ if $\text{prov}(K_i) \mathbin{\&} \text{prov}(K_j) \neq 0$ and they share parameters
3. Topologically sorts via Kahn's algorithm
4. Assigns streams via greedy coloring (lowest stream number without unsatisfied dependency)
5. Emits structured metadata as PTX comments:

\begin{lstlisting}[basicstyle=\ttfamily\small]
// SOUNIO_DAG stream_count=3
// SOUNIO_STREAM kernel=0 stream=0
// SOUNIO_STREAM kernel=1 stream=1
// SOUNIO_DEP from=0 to=2
\end{lstlisting}

A host glue emitter reads this metadata at compile time and generates multi-stream launch code with \texttt{cudaStreamCreate}, per-kernel stream assignment, and dependency-based \texttt{cudaStreamSynchronize} insertion.

\subsection{4 Speculative Epistemic Execution}

When a kernel pipeline contains both coarse and refinement stages, the compiler classifies kernels by their computational cost (presence of transcendental operations, iteration depth) and inserts guards:

\begin{lstlisting}[basicstyle=\ttfamily\small]
if aggregate_epsilon > threshold:
    launch_refinement_kernel(...)  // expensive
// else: coarse result is sufficient
\end{lstlisting}

The guard metadata (\texttt{SOUNIO_SPEC_GUARD}) is embedded in the PTX and interpreted by the host glue, enabling runtime decisions without recompilation.

\subsection{5 Epistemic Kernel Fusion}

Standard kernel fusion scores candidate pairs by parameter overlap and register pressure. Our fusion scoring adds a \emph{provenance diversity bonus}:

\[\text{score} = \text{base\_score} + 60 \times \text{popcount}(\text{prov}_A \oplus \text{prov}_B)\]

Kernels with higher provenance diversity (many differing bits in the XOR) are favored as fusion candidates, as a heuristic for reduced source overlap. We note that provenance diversity does not prove statistical independence --- different origin bits suggest low overlap of measurement classes, but do not guarantee uncorrelated uncertainties. A more interpretable metric (e.g., Jaccard distance between bitsets) could normalize by union size; we use popcount for simplicity.


\section{Evaluation}

\subsection{1 Programmer Effort}

We compare the code required to compute drug concentration with GUM uncertainty from four sensors (blood draw, CT scan, MRI, in-vitro assay):

\begin{table}[t]\centering
\begin{tabular}{lll}
\toprule
Component & CUDA C++ & Sounio \\
\midrule
Dual struct (\texttt{value}, \texttt{variance}) & 10 lines & 0 (native type) \\
Propagation helpers (add, mul, div, sqrt) & 60 lines & 0 (compiler-generated) \\
Shadow array allocation & 25 lines & 0 (automatic) \\
Modified kernel body & 40 lines & 0 (same kernel) \\
Host-side GUM budget & 30 lines & 0 (built-in) \\
Provenance tracking & 25 lines & 0 (automatic) \\
\textbf{Total additional effort} & \textbf{~190 lines} & \textbf{0 lines} \\
\bottomrule
\end{tabular}
\end{table}

\subsection{2 Overhead Comparison}

\begin{table}[t]\centering
\begin{tabular}{lllll}
\toprule
Method & Overhead factor & GPU? & Analytical? & Source \\
\midrule
Sounio shadow lanes & 4–10$\times$ instr. count (exact model); 2.19$\times$ L4 GEMM latency (preliminary) & Yes & Yes (GUM) & Sections 3.3, 5.2 \\
Manual CUDA C++ & Same (if correct) & Yes & Yes & Manual \\
Python \texttt{uncertainties} \cite{uncertainties} & ~1,400x (reported) & No & Yes & Issue tracker \cite{uncertainties_issue57} \\
Julia \texttt{Measurements.jl} \cite{measurementsjl} & ~50–1,500x (reported) & No & Yes & Issue tracker \cite{measurements_issue25} \\
Monte Carlo ($10^4$ samples) & $\sim$10,000x & Yes & No & By construction \\
Monte Carlo ($10^6$ samples) & $\sim$1,000,000x & Yes & No & By construction \\
\bottomrule
\end{tabular}
\end{table}

\textbf{Important caveats.} The instruction-count ratio in Section 3.3 is an exact arithmetic model, not a measured wall-clock bound. Preliminary NVIDIA L4 (Ada Lovelace, sm\_89) GEMM measurements show a 2.19$\times$ geometric-mean latency ratio versus cuBLAS SGEMM at dimensions 1024–8192; these runs used \texttt{epistemic_enabled=false} and therefore do not yet isolate shadow-register overhead from compiler maturity effects. Full shadow-lane wall-clock characterization remains future work (Section 7.3). The overheads reported for \texttt{uncertainties} and \texttt{Measurements.jl} come from user-filed issue reports, not controlled benchmarks; we cite them as motivation, not as rigorous comparisons. Unlike Monte Carlo, analytical propagation provides exact first-order results in a single pass.

\subsection{3 Correctness}

The implementation passes 22 test cases verifying:

\begin{itemize}
\item GUM algebraic properties: commutativity, scaling identity, quadrature, non-negativity
\item Provenance properties: OR idempotency, commutativity, associativity, bit-tagging
\item Sensor fusion: weighted mean uncertainty reduction below best individual source
\item Entropy dispatch: correct variant selection for concentrated vs. spread distributions
\item Per-module self-tests: 10/10 for each novelty module (fusion, speculative, DAG, warp-vote, entropy)
\end{itemize}

\subsection{4 Implementation Scale}

The epistemic GPU stack comprises 9,122 lines of Sounio across 10 self-hosted modules, including the shadow lane emitter, five optimization passes, and three host glue generators (PTX, Metal, SPIR-V). The compiler is fully self-hosted.


\section{Related Work}

\textbf{Table 1. Design-space positioning.} Rows summarize what each system propagates, whether it targets GPU kernels, and whether it is compiler-integrated. “Analytical GUM” means first-order variance quadrature per JCGM 100:2008, not Monte Carlo or full distribution grids.

\begin{table}[t]\centering
\begin{tabular}{llllll}
\toprule
System & Quantity propagated & GPU kernels & Compiler-integrated & Analytical GUM & WMMA shadows \\
\midrule
Puffin \cite{puffin2021} & Epistemic + aleatory UQ & No & Yes (source rewrite) & Intrusive & No \\
Enzyme \cite{enzyme2021} & Jacobian entries & Yes & Yes (LLVM IR) & No (AD chain rule) & No \\
Uncertain$\langle T \rangle$ \cite{uncertaint2014} & Monte Carlo samples & No & No (library type) & No & No \\
\texttt{uncertainties} \cite{uncertainties} / \texttt{Measurements.jl} \cite{measurementsjl} & GUM variance & No & No & Yes & No \\
GBEES-GPU \cite{gbeesgpu2025} & Phase-space distributions & Yes & No & No (grid Bayes) & No \\
Fasi et al. \cite{fasi2021} / Khattak & Mikaitis \cite{khattak2025} & Rounding / hardware error & N/A & No & No & Analysis only \\
\textbf{Sounio shadow lanes} & \textbf{GUM variance + provenance} & \textbf{Yes} & \textbf{Yes (self-hosted)} & \textbf{Yes} & \textbf{Yes} \\
\bottomrule
\end{tabular}
\end{table}

\textbf{Automatic differentiation on GPU.} Enzyme \cite{enzyme2021} is the closest structural precedent: shadow registers in LLVM IR for reverse-mode AD through GPU kernels. Our lanes differ in \emph{what} they propagate (combined variance $u_c^2$, not Jacobian entries) and \emph{how} they compose (quadrature sum, not chain rule). The techniques are complementary --- Enzyme could supply sensitivity coefficients for correlated GUM inputs.

\textbf{Uncertainty compilers and CPU libraries.} Puffin \cite{puffin2021} is, to our knowledge, the only prior work explicitly framed as an “uncertainty compiler,” but it targets CPU Python source and intrusive UQ library calls --- not PTX/WMMA emission. The GUM Tree Calculator \cite{gtc}, Python \texttt{uncertainties} \cite{uncertainties}, and Julia \texttt{Measurements.jl} \cite{measurementsjl} provide analytical GUM on CPU with reported overheads of 50–1,500$\times$ [6,7].

\textbf{GPU uncertainty propagation.} GBEES-GPU \cite{gbeesgpu2025} accelerates grid-based Bayesian estimation for nonlinear dynamics on GPU. It evolves full probability distributions over phase space; our shadows track per-register GUM variance alongside ordinary kernel arithmetic.

\textbf{Uncertain types and probabilistic programming.} Uncertain$\langle T \rangle$ \cite{uncertaint2014} demonstrated first-order uncertain types in managed languages via Monte Carlo sampling, with no GPU backend. Pyro \cite{pyro2019} and Stan \cite{stan2017} model uncertainty through Bayesian inference --- inference frameworks, not general-purpose kernel compilers.

\textbf{Tensor core numerics.} Blanchard et al. \cite{blanchard2020}, Fasi et al. \cite{fasi2021}, and Khattak and Mikaitis \cite{khattak2025} characterize rounding behaviour and hardware non-IEEE features of NVIDIA tensor cores. We propagate \emph{measurement uncertainty} attached to input data through WMMA --- a distinct quantity from hardware rounding error.

\textbf{GPU kernel dispatch.} Stream-K++ \cite{streamkpp2024} and Triton's autotuner \cite{triton2019} select kernel variants from performance characteristics. LithOS \cite{lithos2025} atomizes kernels for fine-grained scheduling. None use information-theoretic measures of input uncertainty for dispatch decisions.


\section{Resolved Limitations}

Three structural limitations identified in earlier versions of this work have been addressed (second-order GUM, covariance shadow, divergence cost model). Three additional bugs identified by peer review have since been resolved (see Section 7.1):

\textbf{Second-order GUM propagation (resolved).} The shadow lane emitter now supports optional Hessian correction terms. For $y = f(x)$, the second-order variance is:

\[u^2(y) = (f'(x))^2 \cdot u^2(x) + \tfrac{1}{2} (f''(x))^2 \cdot u^4(x)\]

A curvature threshold gates activation: the correction is applied only when $|f''(x)| \cdot u^2(x)$ exceeds a configurable tolerance, avoiding unnecessary computation for near-linear operations. For the benchmark $y = x^2$ at $x = 3.0, u(x) = 0.1$: the first-order variance is 0.36, the Hessian correction is 0.0002 (0.056% improvement), and the combined variance is 0.3602 --- matching the analytical second-order Taylor expansion.

\textbf{Assumption.} The second-order correction follows GUM Section E.3.2 and is valid under the assumption that input uncertainties are approximately normally distributed. For non-Gaussian inputs (highly skewed, multimodal, or heavy-tailed distributions), the Hessian term provides a curvature-aware \emph{approximation} but not a rigorous statistical bound. Users with non-Gaussian measurement models should verify distributional assumptions or fall back to Monte Carlo propagation per GUM Supplement 1 (JCGM 101:2008) \cite{jcgm101}.

\textbf{Covariance matrix GPU propagation (resolved).} The compiler now supports correlated inputs via an upper-triangular covariance shadow structure in shared memory. For $N$ correlated variables, the full GUM Eq. 13 propagation is:

\[u^2(y) = \sum_i \left(\frac{\partial f}{\partial x_i}\right)^2 u^2(x_i) + 2 \sum_{i < j} \frac{\partial f}{\partial x_i} \frac{\partial f}{\partial x_j} u(x_i, x_j)\]

Storage is upper-triangular: $N(N+1)/2$ entries, capped at $N = 8$ correlated input sources (36 entries, 288 bytes in f64) to fit within GPU shared memory alongside kernel data. Here $N$ is the number of correlated \emph{measurement sources} contributing to a single datum (e.g., $N = 4$ for four sensors measuring the same analyte), not the number of GPU registers or array elements. The cross-term contribution $2 j_0 j_1 \sigma_{01}$ is computed per-element and added to the diagonal-only result.

\textbf{Warp divergence cost model (resolved).} An analytical cost model quantifies the penalty of the dual-path warp-vote mechanism. The model comprises three components:

\begin{itemize}
\item \emph{Serialization penalty}: $P_s = 1.0 + d/100$ where $d$ is the divergence percentage (0% $\rightarrow$ 1.0x, 100% $\rightarrow$ 2.0x)
\item \emph{Ballot overhead}: $P_b = n_b \cdot c_b / C_k$ where $n_b$ = number of ballots, $c_b$ = ballot cost in cycles (~2 on SM 8.x), $C_k$ = typical kernel cycles
\item \emph{Reconvergence cost}: $P_r = L \cdot d/100 \cdot 0.01$ where $L$ = instruction distance between diverge and reconverge points
\end{itemize}

The combined penalty $P = P_s + P_b + P_r$ feeds into a speedup estimator. With budget-bounded fast-path overhead of 0.3x (30% of full shadow cost) and 10% divergence, the predicted speedup is approximately 1.5x. The model populates the \texttt{estimated_speedup} field in \texttt{WarpVotePlan}, enabling cost-aware decisions about whether to apply the dual-path optimization. We emphasize that these are \emph{model predictions} based on instruction counting; actual speedup depends on register pressure, occupancy, and memory access patterns that require hardware measurement.

\subsection{1 Resolved Since Initial Preprint}

Three additional items identified in the initial preprint have now been resolved:

\textbf{Sqrt domain validity (resolved).} \texttt{knowledge_sqrt_f64} now emits a \texttt{setp.ge.f64 %dp, %av, 0d0000000000000000} domain check and folds it into the validity predicate via \texttt{and.pred %lp, %ap, %dp}. Input values $a < 0$ set \texttt{%lp = false} before the sqrt executes; the \texttt{sqrt.rn.f64} instruction may produce NaN but the false validity predicate suppresses downstream propagation. Epsilon propagation uses $u(\sqrt{a}) = u(a)/(2\sqrt{a})$ with a near-zero guard: when $\sqrt{a} = 0$ exactly, \texttt{%le} receives \texttt{%ae} as a conservative upper bound to avoid division by zero.

\textbf{Division domain validity (resolved).} The near-zero denominator guard now also \emph{invalidates} the result, not merely clips the epsilon. The original code emitted \texttt{@%tp mov.f64 %le, 1.0} (cap epsilon) then propagated inherited validity only. The fix adds \texttt{not.pred %tp, %tp; and.pred %lp, %lp, %tp} after the clip: results with near-zero denominators are both clamped and marked invalid.

\textbf{Entropy dispatch scale-dependence (resolved).} \texttt{ed_compute_histogram_normalized} normalizes epsilon values to relative uncertainty $\epsilon / |x|$ before histogramming. When $|x| < 10^{-15}$, raw epsilon is used as a conservative fallback. This makes Shannon entropy $H(\epsilon/|x|)$ dimensionless and independent of input scale or physical units. Callers using physical quantities (e.g., milligrams, meters) should use the normalized variant.

\subsection{2 Remaining Open Limitations}

\textbf{Logarithm domain validity.} \texttt{ln(a)} requires $a > 0$; the current emitter has no \texttt{knowledge_log_f64} intrinsic and no domain check. Future work.

\textbf{Tiled covariance empirical characterization.} The \texttt{tc_propagate_tiled} function (N = 128, CUTLASS-style blocking) compiles and produces correct variance estimates against the N $\leq$ 8 monolithic reference, but has not been benchmarked for shared-memory pressure or occupancy on real hardware.

\subsection{3 Remaining Future Work}

\textbf{Hardware benchmarks.} This preprint reports structural correctness and programmer effort reduction. Wall-clock benchmarks on NVIDIA Ampere (A5000, sm_86) and Ada Lovelace (L4, RTX 4000 Ada, sm_89) hardware are in progress and will be reported in the full paper.

\textbf{Higher-order covariance (implemented, pending empirical characterization).} Tiled covariance propagation extends support to $N = 128$ correlated variables via CUTLASS-style upper-triangular tile blocking. The $N \times N$ covariance matrix is partitioned into $T \times T$ tiles (where $T$ auto-tunes per SM architecture: 64 on Hopper sm\_90, 48 on Ada Lovelace sm\_89, 32 on Ampere sm\_86). Each tile fits in shared memory independently; results accumulate via tree reduction. For $N = 64$: 3 upper-triangular tiles at 4,224 bytes each. For $N = 128$: 10 tiles. The per-tile J$\cdot\Sigma_{\text{tile}}\cdot$J$^\top$ computation reuses the same propagation kernel as the monolithic $N \leq 8$ case. \textbf{Clarification:} $N$ here is the number of \emph{correlated input uncertainty sources} per datum (e.g., $N = 4$ for blood draw + CT + MRI + assay), not the number of GPU registers or kernel dimensions. The covariance matrix is carried per logical measurement group, not per thread.


\section{Conclusion}

Epistemic shadow lanes demonstrate that first-order GUM uncertainty propagation through GPU kernels is a compiler problem, not a library problem. The transformation removes manual per-operation uncertainty propagation and metadata plumbing from GPU kernels; the instruction-count overhead is bounded by a moderate constant factor in the current model, though hardware characterization remains future work. The five optimizations we introduce --- warp-vote fast-paths, entropy-gated dispatch, provenance-aware scheduling, speculative execution, and epistemic fusion scoring --- show that uncertainty metadata is not merely overhead to be tolerated, but information that can guide execution-time decisions in ways unavailable to conventional GPU compilers.


\section{Software Availability and Reproducibility}

Source code, compiler tests, and benchmark artefacts are maintained in the public Sounio repository: <https://github.com/Sounio-lang/sounio>. The epistemic GPU stack lives under \texttt{self-hosted/gpu/}; shadow-lane correctness tests are in \texttt{tests/gpu/}. Preliminary L4 GEMM measurements cited in Section 5.2 are documented in \texttt{benchmarks/results/NVIDIA_L4_BENCHMARKS.md} with machine-readable data in \texttt{benchmarks/results/l4_raw_data.json}. To reproduce the compiler check surface:

\begin{lstlisting}[basicstyle=\ttfamily\small]
export SOUNIO_STDLIB_PATH=$(pwd)/stdlib
./bin/souc check self-hosted/gpu/hlir_to_gpu.sio
./bin/souc run tests/run-pass/epistemic_kernel_shadow.sio
bash scripts/run_sio_test_suite.sh epistemic_kernel_shadow
\end{lstlisting}

GPU wall-clock reproduction requires NVIDIA CUDA hardware and the dispatch scripts referenced in the benchmark report.


\bibliographystyle{plain}
\bibliography{refs}
\end{document}
